% ============================================================
% ［架空の記入例・提出不可］結果，解釈および今後の課題はすべて架空である．
% 通常のchapters/06-conclusion.texには自動で読み込まれない．
% ============================================================
\chapter{\UsamiConclusionTitle}

\begin{center}
  \UsamiFictionalExampleNotice
\end{center}

\section{本研究のまとめ}

本研究では，少数症例下の腹部CT多臓器分割を対象として，解剖構造を保持する条件付き潜在拡散データ拡張を検討した．臓器ラベルを条件として合成CTを生成し，臓器体積，左右臓器の位置関係および臓器間重複を検査した上で，採択された合成対だけを3D U-Netの学習へ追加した．また，分割学習の目的関数へ解剖整合性項を導入し，生成段階の品質管理と分割段階の構造制約を組み合わせた．

架空の患者単位評価では，提案法の臓器平均Dice係数は$0.826$，HD95は\SI{9.7}{mm}であり，実画像のみの条件に対してDiceが0.045高く，HD95が\SI{3.1}{mm}低かった．一方，細い血管領域では条件間の差が小さく，すべての臓器で一様な改善は得られなかった．解剖整合性項を除いた架空の条件では平均Diceが0.011低下したが，この差だけから生成画像の臨床的妥当性や外部施設での再現性を結論づけることはできない．

以上より，本例の条件では，合成対を無条件に追加するのではなく，ラベルとの対応と解剖構造を検査してから分割学習へ利用する設計が有効である可能性を示した．本章の症例数，性能値および解釈は，修士論文のための架空例であり，実在する研究成果ではない．

\section{研究の限界}

第1に，単一施設・単一造影相を想定しており，撮像装置，造影タイミングおよび再構成条件の変化を評価していない．第2に，臓器体積と位置関係による採択は粗い品質管理であり，局所境界の不整合や病変による正常解剖からの逸脱を適切に扱えない可能性がある．第3に，分割性能をDiceとHD95で評価したが，放射線画像に通じた評価者による合成画像の盲検判定，臨床ワークフローでの有用性および生成モデルのプライバシー監査は実施していない．

\section{今後の課題}

今後は，施設と造影相を分けた外部評価を行い，症例分布が変化した場合のDice，HD95および校正誤差を確認する必要がある．次に，臓器間の接触面や血管分岐を扱える構造表現を導入し，採択判定の偽陽性と偽陰性を評価する．さらに，合成対の追加枚数を揃えた対照実験と患者単位の信頼区間を用いて，単なるデータ量の増加と解剖構造保持の効果を分離する．

実データへ適用する場合は，倫理審査，匿名化，利用範囲および生成データの保管方針を確定した後，学習症例の再同定可能性と暗記を監査する．これらを満たして初めて，データ拡張手法としての一般化可能性を判断できる．
